\documentclass[11pt]{article}
\usepackage{amsmath,amssymb,amsthm,array}
\usepackage[margin=28mm]{geometry}
\setlength{\emergencystretch}{3em}
\newtheorem{theorem}{Theorem}
\newtheorem{lemma}[theorem]{Lemma}
\newtheorem{corollary}[theorem]{Corollary}
\newcommand{\Fix}{\operatorname{Fix}}
\newcommand{\Stab}{\operatorname{Stab}}
\begin{document}
\title{Joint fixed sets and the remaining original prefix}
\author{Human Compression}
\date{}
\maketitle

\begin{abstract}
Two descriptions of the same joint fixed set can force every member
of a string marking to preserve it. This gives an obstruction to
generation that uses only the original permutation action. We prove
the criterion and specify its application to the exact prefix of
original job $[68,5]$. The complete exploratory fixed sets agree and
suggest a nine-point certificate. That exploration is retained as a
diagnostic; a separately qualified producer and strict Saved check,
explicit reuse of the target group's original generator lookup, and independent
actual review remain required before job adoption.
\end{abstract}

\section{The set criterion}

Let an ambient group $\Gamma$ act on a set $\Omega$, and let
$G\leq\Gamma$ be the target group, with right-action notation
$x\mapsto x^g$. For $g\in\Gamma$, write
$\Fix(g)=\{x\in\Omega:x^g=x\}$. A string marking of rank $r$
is a sequence $(t_0,\ldots,t_{r-1})$ satisfying
\[
                  [t_i,t_j]=1\quad\text{if }|i-j|\geq2.
\]
In the usual applications the marks are involutions. The following
criterion does not require that additional assumption.

\begin{lemma}\label{lem:centralizer}
If $g$ commutes with $u$ and $v$, then $g$ preserves
$\Fix(u)\cap\Fix(v)$ setwise.
\end{lemma}
\begin{proof}
For $x$ fixed by $u$ and $v$, commutation gives
\[
        (x^g)^u=(x^u)^g=x^g,\qquad
        (x^g)^v=(x^v)^g=x^g.
\]
Thus the image of the joint fixed set under $g$ is contained in
that set. The same argument applies to $g^{-1}$ and gives the reverse
inclusion. This works for infinite as well as finite sets.
\end{proof}

\begin{theorem}\label{thm:gluing}
Let the marks belong to $\Gamma$, and suppose $r\geq5$.
Put $a=t_0$, $b=t_1$ and $e=t_{r-1}$.
If
\begin{equation}\label{eq:sets}
 S=\Fix(a)\cap\Fix(e)=\Fix(a)\cap\Fix(b),
\end{equation}
then every mark preserves $S$ setwise. If there are an actual member
$g\in G$ and a point $x\in S$ with $x^g\notin S$, then
$H=\langle t_0,\ldots,t_{r-1}\rangle\ne G$. If the marks also
belong to $G$, then $H<G$.
\end{theorem}
\begin{proof}
The two definitions of $S$ show that $a,b,e$ each fix it pointwise.
For every index $i$ with $2\leq i\leq r-3$, the string relations
make $t_i$ commute with both $a$ and $e$. These marks preserve $S$
by Lemma~\ref{lem:centralizer}, using the first description in
\eqref{eq:sets}.

The only remaining index is $i=r-2$. Since $r\geq5$, its distance
from indices $0$ and $1$ is at least two. The penultimate mark
therefore commutes with $a$ and $b$, so it preserves the second
description of $S$. This accounts for every mark, including the
boundary case $r=5$, where the middle range contains only $i=2$.

The subgroup they generate lies in $\Stab_\Gamma(S)$. The displayed
escape gives $g\notin H$, whereas $g\in G$. Consequently $H\ne G$,
regardless of whether the marks belong to the target group. If their
membership in $G$ is separately known, then $H\leq G$ and hence
$H<G$.
\end{proof}

We may take $\Gamma=\operatorname{Sym}(\Omega)$, so no original
membership proof for the prescribed permutations is needed for the
negative generation conclusion. The escape itself implies that
$S$ is nonempty and proper. Neither
faithfulness nor transitivity of the action is used. No field,
linear representation, characteristic-vector computation, order
calculation or intersection property is needed. In particular,
the theorem does not require $a$ and $b$ to commute.

The setwise-stabilizer escape principle is already used in the
project's \emph{Original invariant sets}, edition~2,
\texttt{original\_invariant\_sets/v2/ARGUMENT.rst}. The new
specialization here is the use of the two centralizer descriptions
in \eqref{eq:sets} to cover every position of a string marking.
Root proposed this equality criterion after the first branch~68
exploration. No claim of historical priority is made. The separate
characteristic-vector and maximal-endpoint theorems remain useful
other criteria; no old operational verdict is transferred to this
new predicate.

\section{The exact original prefix}

The original job is $[68,5]$. The exact whole pair file is
\texttt{pair05.g}, 549385 bytes, with SHA256
\begin{center}\small
\texttt{594ce69d6dd535c6178f87c9515805312fa6f0d69bb6a9bd89c48142643982ca}.
\end{center}
Its ordered entries are $(a,e)$. The exact whole third-mark file is
\texttt{branch0068.g}, 276591 bytes, with SHA256
\begin{center}\small
\texttt{7aa368f3b42f704c1ea1e2047dc7844dd3215f78327584b547113ba994761723}.
\end{center}
Its value is $b$. The original end-to-end traversal stores the
three values as $(a,e,b)$ but places them at marking positions
$(t_0,t_{r-1},t_1)$. Thus the rank-five prefix is
\[
                              (a,b,\_,\_,e).
\]
This order is essential: neither a source ordinal nor a third mark
from another branch may replace this literal original $b$.

The target original group/model premise is the whole
\texttt{q8\_model.g}, 733634 bytes, SHA256
\begin{center}\small
\texttt{5d08810d95f0c9454fbbd687133a23634e3a1c30c78b9c71cb4ea370cfa3cd0a}.
\end{center}
Its named-group and software meanings remain imported premises. The
model defines the target variable $G$ before the separate variable
$A$, with model name \texttt{G2\_8}. The target here is that original
$G$, not $A$, the full automorphism group or its symplectic embedding.
The accepted endpoint reversal replay
\texttt{acd78878\ldots 0583c8d} supplies the exact original
pair-membership limb for pair~05. The whole source, receipt, output
and model joins of that acquired replay remain unchanged.

There is also an existing positive original-membership limb for this
exact $b$. The native source
\texttt{q8\_high\_adaptive2\_branch0068.g}, 1987 bytes,
SHA256
\begin{center}\small
\texttt{57c508ee364c3c6761ca7ca92c756ee3b5a328dd5773279a3b03b18e0845b2c7},
\end{center}
reads this whole pair and third mark. Before its later search, it
explicitly rejects any of $a,e,b$ outside the same original $G$ or
not of order two. Its receipt, 3274 bytes,
\texttt{14022016\ldots 4047ac6}, records return code zero, the
complete marker, empty standard error and unchanged inputs. Its
complete output is 320 bytes,
\texttt{91fa7c40\ldots a546ce}.

This positive limb remains attributed historical context. The ambient
form of Theorem~\ref{thm:gluing} makes it unnecessary for the new
negative argument. No new membership gate, native replay or pool
import is required. A positive membership claim made for another
purpose would require an explicit reuse acceptance retaining the
whole subjects and software qualifications. The old empty leaf list
and normalizer-pruning output are not used as a negative conclusion.
The new scientific check must
independently verify the whole permutation identities
\[
 a^2=b^2=e^2=1,\qquad [a,e]=[b,e]=1.
\]
It must \emph{not} impose $[a,b]=1$: the historical first bond
has order four. These finite identities establish the prescribed
involutory string prefix. They do not assert target-group membership,
which is unnecessary for the exclusion.

\section{The proposed finite certificate}

The current exploration reports the complete sets
\begin{equation}\label{eq:nine}
 \Fix(a)\cap\Fix(e)=\Fix(a)\cap\Fix(b)
       =\{1,10,11,12,13,14,15,16,17\}.
\end{equation}
It also reports the same triple intersection, with individual fixed
counts $585,73,73$ for $a,e,b$. The first earlier scan established
only the $a,b$ set. The equality in \eqref{eq:nine} cannot be
deduced from two equal cardinalities. It is the equality of the two
complete sorted original point lists that is required.

The exploratory output has zero Source, proof and job credit. The
qualified producer and strict Saved relation are to reconstruct both
complete sets from the whole original inputs, rather than import
\eqref{eq:nine} as an accepted finite result. Their certificate
contains all three vector digests, all three fixed counts, the full
$AE$, $AB$ and triple lists, all required involution/commutation
checks, and the supplied escape. A valid but inapplicable outcome
must remain unpaid. Saved checks the supplied certificate directly;
it does not call the producer or search for an escape.

An already acquired original ambient lookup is
\begin{equation}\label{eq:lookup}
            10^{g_2}=586,
\end{equation}
where $g_2$ is original generator~2 of the target variable $G$ in
this same model. It is not a generator of $A$ or an unspecified
ambient automorphism. Both generator
and point labels are \emph{one-based}. This edge was actually checked
by the preceding strict Saved profile computation, whose whole
4522-byte certificate has SHA256
\begin{center}\small
\texttt{451c6dec5769465cfee483581f070e4a0084eb5ccb40dca32afe1b65f544768c}.
\end{center}
Its original generator meaning is bound to the accepted complete
projective bridge, whose certificates have SHA256
\texttt{0601a736\ldots 184b4f8}, and to the whole bridge gate
\texttt{4ecaf56b\ldots 3802648}. The current Root application
acceptance \texttt{2ba1396a\ldots e6c4aac3} binds that actual
history and records the original one-based edge.

The bridge field \texttt{ordered\_original\_generator\_images}
retains the ordered original $G$ generators. Its second entry is
the target member used here. The lookup value \eqref{eq:lookup}
is independent of the endpoint pair: it is a value of that same
original target permutation.
Consequently it may be reused as an attributed lookup premise.
The pair~08 fixed-set predicate from the preceding result is not
reused for branch~68. The new finite check must establish that $10$
belongs to the new common set and that $586$ does not. A Root reuse
gate must bind the exact original model, target $G$ generator index,
lookup, whole accepted bridge and whole producer/Saved history.
No new matrix recovery or ambient-generator computation is needed.

\begin{corollary}\label{cor:branch68}
Assume the exact original prefix-file/frontier ordering and target
original $G$ model/lookup premises above.
If a qualified complete producer/Saved certificate establishes
\eqref{eq:nine}, and the original lookup \eqref{eq:lookup} is
explicitly authenticated and reused, then no rank-five completion of
original branch~68 generates $G$.
\end{corollary}
\begin{proof}
The exact ordering of the original prefix makes $a,b,e$ occupy the
positions required by Theorem~\ref{thm:gluing}. Equality of the
complete fixed sets gives its set $S$. The point $10$ belongs to
the displayed nine-set and $586$ does not. The authenticated lookup
therefore supplies an actual member $g_2\in G$ moving a point of
$S$ outside $S$. The theorem proves the claim.
\end{proof}

This is a separate set argument for a fixed original prefix. It is
not an application of the maximal-endpoint theorem to pair~05.
That endpoint profile remains $(585,73,9)$ and outside the two
rank-three-endpoint criterion. The present implication needs the
actual original $b$, complete fixed-set equality and ambient escape.

\section{Scope of any subsequent adoption}

The current fixed required universe has 686 jobs. After the separately
completed 76-job registration it has 685 closed jobs and the single
pending job $[68,5]$. The corollary would justify removing that
one job only after the new Source, runtime, controls, complete
producer/Saved, distinct actual review and typed Root application
requirements have been met. Those are future operations; this
document changes no state.

All old 685 proofs, all rank-four records, the original frontier and
software qualifications, both protected progress files and unrelated
state must remain intact. Closing the last job in this fixed universe
would still not by itself establish global frontier exhaustion, the
full automorphism/reversal census, a whole original-manuscript
recovery, a conjecture result, numerical compression or Lean
formalization. No such credit is claimed here.
\end{document}
